\documentclass[11pt,a4paper]{article}

% ── Packages ──────────────────────────────────────────────────────────
\usepackage[utf8]{inputenc}
\usepackage[T1]{fontenc}
\usepackage{lmodern}
\usepackage[margin=1in]{geometry}
\usepackage{amsmath,amssymb,amsthm}
\usepackage{booktabs}
\usepackage{graphicx}
\usepackage{hyperref}
\usepackage{algorithm}
\usepackage{algpseudocode}
\usepackage{xcolor}
\usepackage{enumitem}
\usepackage{subcaption}
\usepackage{bm}

% ── Macros ────────────────────────────────────────────────────────────
\newcommand{\R}{\mathbb{R}}
\newcommand{\E}{\mathbb{E}}
\newcommand{\argmin}{\operatorname*{arg\,min}}
\newcommand{\argmax}{\operatorname*{arg\,max}}
\DeclareMathOperator*{\diag}{diag}

% ── Title ─────────────────────────────────────────────────────────────
\title{Vector Policy Optimization:\\Learning from Structured Privileged Feedback\\[6pt]
{\large Draft --- For Internal Discussion}}
\author{Terminal Contrastive RL Project}
\date{June 2026}

\begin{document}
\maketitle

\begin{abstract}
Standard reinforcement learning from human feedback (RLHF) collapses all
feedback into a single scalar reward.  When a \emph{privileged teacher}
can provide structured, multi-axis feedback---answer correctness,
reasoning consistency, tool-call validity, brevity---scalarisation
discards precisely the information the teacher was introduced to inject.
We propose \textbf{Vector Policy Optimization (VPO)}, a family of
update rules that preserve the vectoriality of the reward signal at the
gradient level.  We present a spectrum from fixed-weight scalarisation
through Dirichlet random mixing (already validated in our Vector-$\lambda$
V2 experiments, $+16.1$ pp over base) to \textbf{MGDA-VPO}, which
computes per-axis policy gradients and combines them via a convex
optimisation on the simplex that guarantees Pareto-stationary updates.
We sketch a concrete experiment (Exp19), analyse when VPO earns its keep
vs.\ degenerating to averaging, and identify open questions around critic
architecture, compute cost, and stability.
\end{abstract}

\tableofcontents
\newpage

% ======================================================================
\section{Motivation: Why Scalar Rewards Waste Privileged Information}
\label{sec:motivation}

The standard policy gradient maximises expected advantage:
\begin{equation}
  \nabla J(\theta) \;=\; \E_{\tau \sim \pi_\theta}\!\Big[
    \nabla \log \pi_\theta(\tau) \;\cdot\; A(\tau)
  \Big],
  \qquad A(\tau) \in \R.
  \label{eq:pg-scalar}
\end{equation}
The advantage $A(\tau)$ is a single number.  When a privileged teacher
returns a vector of feedback signals $\bm{R}(\tau) \in \R^K$, the
conventional pipeline reduces it to a scalar via a fixed or random
weighted sum \emph{before} computing the gradient.  This is informationally
lossy: a trajectory that is right on $K{-}1$ axes but wrong on one gets
the same scalar as a trajectory that is mediocre on all axes.

\paragraph{Example.}
Consider a bash agent rollout that produces the correct answer but uses
an invalid tool-call format.  A 5-axis teacher returns
\[
  \bm{A}(\tau) = (\underbrace{+1}_{\text{correct}},\;
                   \underbrace{+1}_{\text{reasoning}},\;
                   \underbrace{-1}_{\text{tool format}},\;
                   \underbrace{+1}_{\text{state match}},\;
                   \underbrace{-0.5}_{\text{verbose}}).
\]
Scalarising with equal weights gives $A = +0.3$---``this trajectory was
$0.3$ good.''  The gradient says ``do more of this, a bit,'' and the
model slowly drifts.  A vector update can say: ``fix the tool call; the
rest is fine.''  That targeted signal is exactly what pedagogical RL is
for, and scalarisation destroys it.

\paragraph{The exploration problem.}
On-policy RL already struggles when the policy rarely samples correct
trajectories: the gradient has no positive examples to learn from.
Pedagogical RL addresses this by injecting trajectories or rewards the
policy could not produce itself.  Scalarising the teacher's feedback
partially negates this injection: we brought in structured information
and then threw away the structure.

% ======================================================================
\section{Background: What We Already Know Works}
\label{sec:background}

\subsection{Scalar self-similarity (Exp17)}

The simplest verifier-free reward: for each rollout $i$ in a group of
$G{=}8$, compute mean pairwise terminal similarity
\[
  z_i = \frac{1}{G{-}1}\sum_{j \neq i} \mathrm{sim}(\mathrm{out}_i,\,
        \mathrm{out}_j),
\]
then z-score over the group and use as the GRPO advantage.
Result: $+12.5$ pp over base on Prime A6000 ($+13.4$ pp at step 150).

\subsection{Variational V1 (Exp9)}

Same scalar $z_i$ but with an explicit variational interpretation and
weight sync.  Result: $+15.2$ pp over base on nodeset3.

\subsection{Vector-$\lambda$ V2 (Exp10)}

Three reward axes per rollout:
\begin{equation}
  \bm{R}_i = \big(\mathrm{strict\_sim}_i,\;\;
                \mathrm{jaccard\_sim}_i,\;\;
                \mathrm{exit\_success}_i\big) \in \R^3.
  \label{eq:v2-reward}
\end{equation}
Per training step, sample $\bm{\lambda} \sim \mathrm{Dir}(\alpha{=}1.0)$
and collapse:
\begin{equation}
  R_i = \bm{\lambda} \cdot \bm{R}_i, \qquad
  A_i = \frac{R_i - \bar{R}}{\mathrm{std}(R)}.
  \label{eq:v2-dirichlet}
\end{equation}
Result: $+16.1$ pp over base---the best single number, but only $\sim$1 pp
above V1.  The small gap is consistent with the axes being highly
correlated (both $\mathrm{strict\_sim}$ and $\mathrm{jaccard\_sim}$ measure
output similarity).

% ======================================================================
\section{Spectrum of Vector-Reward Update Rules}
\label{sec:spectrum}

Given a vector advantage $\bm{A}(\tau) \in \R^K$, the design question is:
\emph{how do we turn it into a parameter update?}

\begin{table}[h]
\centering
\caption{Spectrum of vector-reward update rules.}
\label{tab:spectrum}
\begin{tabular}{@{}lll@{}}
\toprule
\textbf{Approach} & \textbf{Update rule} & \textbf{When it earns its keep} \\
\midrule
Fixed linear weight &
  $A_{\mathrm{scalar}} = \bm{w} \cdot \bm{A}$ &
  Baseline only; requires hand-tuning \\[4pt]
\textbf{Dirichlet mixing (V2)} &
  $\bm{\lambda} \sim \mathrm{Dir}(\alpha)$ per step &
  Random regularisation across axes \\[4pt]
Worst-axis (max-min) &
  $A_{\mathrm{scalar}} = \min_k A_k$ &
  Guarantees improvement on every axis; slow \\[4pt]
\textbf{MGDA} (D\'{e}sid\'{e}ri 2012) &
  $\argmin_{\bm{\alpha} \in \Delta^{K-1}} \|\sum_k \alpha_k \bm{g}_k\|^2$ &
  Pareto-stationary; auto-weights complementary axes \\[4pt]
Adaptive per-axis &
  Weight axis $k$ $\propto$ policy weakness on $k$ &
  Spends gradient where needed; needs per-axis eval \\[4pt]
Multi-head policy &
  $K$ heads, each maximises own $A_k$; mix at inference &
  Each axis gets clean signal; mixing rule open \\
\bottomrule
\end{tabular}
\end{table}

\paragraph{Key distinction.}
Approaches 1--3 collapse the vector to a scalar \emph{before} the
gradient; only MGDA and multi-head preserve vectoriality \emph{at the
gradient level}.  This is the crucial difference: if you scalarise
before the backward pass, the gradient has no way to ``unmix'' the
information.

% ======================================================================
\section{MGDA-VPO: Algorithm}
\label{sec:mgda-vpo}

MGDA (Multiple Gradient Descent Algorithm; D\'{e}sid\'{e}ri 2012)
computes per-objective gradients and finds a convex combination that
decreases all objectives simultaneously.  Applied to RL with $K$
advantage axes:

\begin{algorithm}[h]
\caption{MGDA-VPO}\label{alg:mgda-vpo}
\begin{algorithmic}[1]
\Require Policy $\pi_\theta$, $K$ advantage functions $\{A_k\}_{k=1}^K$,
  learning rate $\eta$, group size $G$
\For{each training step}
  \State Generate $G$ rollouts $\{\tau_i\}_{i=1}^G$ via $\pi_\theta$
  \For{$k = 1, \ldots, K$}
    \State Compute per-rollout advantage $A_k(\tau_i)$ over the group
           (z-score normalised)
    \State Compute gradient $\bm{g}_k \gets
           \E\big[\nabla_\theta \log \pi_\theta(\tau) \cdot A_k(\tau)\big]$
           \Comment{$K$ backward passes}
  \EndFor
  \State Solve convex QP on the simplex:
         \[
           \bm{\alpha}^* = \argmin_{\bm{\alpha} \in \Delta^{K-1}}
                           \big\|\textstyle\sum_{k=1}^K \alpha_k \bm{g}_k\big\|^2
         \]
  \State Update: $\theta \gets \theta + \eta \sum_{k=1}^K \alpha_k^* \bm{g}_k$
\EndFor
\end{algorithmic}
\end{algorithm}

\paragraph{Properties.}
\begin{itemize}[nosep]
\item When all $K$ gradients point in a similar direction, $\bm{\alpha}^*$
  is near-uniform and MGDA-VPO reduces to gradient averaging.
\item When axes conflict, $\bm{\alpha}^*$ picks the compromise direction
  that improves \emph{all} axes simultaneously.
\item If no such direction exists (Pareto front reached), $\|\sum
  \alpha_k^* \bm{g}_k\|^2 = 0$ and the update is zero---this is
  \emph{not} a bug but a signal: the policy is Pareto-stationary and
  further improvement requires relaxing one axis.
\end{itemize}

\paragraph{Closed-form for $K \leq 3$.}
For $K = 2$, the QP has a closed-form solution (Sener \& Koltun 2018).
For $K = 3$ (our case with strict/jaccard/exit), the QP can be solved
analytically or via a single call to a tiny QP solver with negligible
overhead.

\paragraph{Compute cost.}
MGDA-VPO requires $K$ backward passes per training step.  For $K = 3$
this is $3\times$ the backward cost.  Mitigation strategies:
\begin{itemize}[nosep]
\item Use a single fused backward with $K$ loss heads (accumulate
  per-axis $\bm{g}_k$ without materialising $K$ separate
  backward-call overhead).
\item Use \emph{ MGDA-Approx}: compute one backward with the
  Dirichlet-mixed advantage (as in V2), then a second with the
  \emph{worst-axis} advantage; combine the two gradients via MGDA.
  This captures most of the benefit at $2\times$ cost.
\end{itemize}

% ======================================================================
\section{When Does VPO Earn Its Keep?}
\label{sec:when}

The central question is not ``is VPO good?'' but rather:

\begin{quote}
\emph{When are the $K$ reward axes carrying independent information
that a scalar cannot represent?}
\end{quote}

\subsection{The degeneracy risk}

If the per-axis advantages are highly correlated estimates of the same
underlying ``goodness''---as is likely for $\mathrm{strict\_sim}$ and
$\mathrm{jaccard\_sim}$, both of which measure output similarity---then
MGDA-VPO degenerates to gradient averaging and we have paid $K\times$
compute for nothing.

Formally: if $A_k \approx c_k \cdot A_0 + \epsilon_k$ with
$\mathrm{Var}(\epsilon_k) \ll \mathrm{Var}(c_k A_0)$ for all $k$,
then $\bm{g}_k \approx c_k \bm{g}_0$ and the simplex solution places
all weight on the dominant direction regardless of $\bm{\alpha}$.

\subsection{Designing uncorrelated axes}

VPO earns its keep when the axes are \emph{deliberately uncorrelated}.
In the privileged-teacher setting, this means choosing feedback
dimensions that probe independent facets of quality:

\begin{table}[h]
\centering
\caption{Candidate reward axes for VPO in the bash-agent domain.}
\label{tab:axes}
\begin{tabular}{@{}llc@{}}
\toprule
\textbf{Axis} & \textbf{Source} & \textbf{Correlation with correctness?} \\
\midrule
Answer correctness & Ground-truth oracle & Perfect (dominant axis) \\
Tool-call format validity & Syntactic check & Moderate \\
Reasoning consistency & LLM judge or heuristic & Moderate \\
Brevity (negative length) & Cheap, mechanical & Low \\
Intermediate env-state match & Privileged env introspection & Low \\
Exit code 0 & Mechanical & Moderate \\
\bottomrule
\end{tabular}
\end{table}

The key insight: axes 3--5 have \textbf{low} correlation with axis 1
(correctness).  A trajectory can be correct but verbose, or correct but
use a different internal state path.  These are exactly the dimensions
where a scalar advantage would conflate ``correct but suboptimal'' with
``partially correct but efficient,'' and where VPO's gradient-level
vector preservation can separate the two signals.

\subsection{Empirical test}

The simplest test for axis independence: after collecting $G{=}8$
rollouts per task across the training set, compute the pairwise Pearson
correlation matrix of the $K$ advantages.  If all off-diagonal entries
$|\rho_{jk}| > 0.8$, VPO will not help.  If some $|\rho_{jk}| < 0.5$,
VPO has room to contribute.

% ======================================================================
\section{Proposed Experiment: Exp19 --- MGDA-VPO}
\label{sec:exp19}

A drop-in replacement for V2's Dirichlet mixing step.

\paragraph{Design.}
\begin{enumerate}
\item Reuse V2's three reward axes: $\mathrm{strict\_sim}$,
  $\mathrm{jaccard\_sim}$, $\mathrm{exit\_success}$.
\item Compute three z-scored advantages $A_{\mathrm{strict}}$,
  $A_{\mathrm{jaccard}}$, $A_{\mathrm{exit}}$ (instead of one
  Dirichlet-mixed advantage).
\item Compute three policy gradients via three masked-backward passes
  (or one fused pass with three loss heads).
\item Solve the 3-simplex MGDA QP analytically per step.
\item Compare to V2 (Dirichlet-mixed) and V1 (single scalar).
\end{enumerate}

\paragraph{Expected outcome with correlated axes.}
Since $\mathrm{strict\_sim}$ and $\mathrm{jaccard\_sim}$ are correlated,
MGDA-VPO will behave like averaging and land within $\pm 2$ pp of V2.
This is the \emph{null result}---valuable because it confirms the
degeneracy hypothesis.

\paragraph{The interesting variant.}
Add deliberately uncorrelated axes---a brevity penalty and a
tool-call-format validity check---and see whether MGDA-VPO opens a gap.
If it does, the contribution is: \emph{VPO with well-designed
uncorrelated axes outperforms any scalarisation.}

\paragraph{Experimental controls.}
\begin{itemize}[nosep]
\item Same environment as Exp9/10 (nodeset3 H100, same bash sandbox).
\item Same $\mathrm{mbs}=8$, $G=8$, $\mathrm{lr}=10^{-5}$,
  $\mathrm{lora\_rank}=16$.
\item Same weight-sync protocol.
\item $\mathrm{max\_steps} = 200$ with $\mathrm{save\_steps} = 50$ for
  early-stopping analysis.
\end{itemize}

% ======================================================================
\section{Relation to Existing Work}
\label{sec:related}

\paragraph{Multi-task learning.}
MGDA was introduced for multi-task learning by Sener \& Koltun (2018),
who showed it produces Pareto-stationary updates for $K$ task losses.
Our contribution is applying it to the \emph{advantage decomposition
within a single RL step}, where each ``task'' is a reward axis from a
privileged teacher.

\paragraph{Multi-objective RL (MORL).}
MORL learns a policy that trades off multiple objectives, typically by
computing a Pareto front over policies.  VPO is different: we do not
seek a set of Pareto-optimal policies but a \emph{single update
direction} that improves all axes simultaneously at each step.

\paragraph{Reward shaping / potential-based shaping.}
Ng et al.\ (1999) showed that potential-based reward shaping does not
change the optimal policy.  VPO is not reward shaping: it changes the
\emph{gradient direction}, not the reward value, and therefore can
change which policy is optimal.

\paragraph{RLHF with process reward models (PRMs).}
PRMs provide dense per-step rewards but still scalarise them.
VPO would apply if the PRM had multiple independent heads (e.g.\
correctness, safety, helpfulness) that were preserved at the gradient
level rather than collapsed into a single score.

% ======================================================================
\section{Open Questions}
\label{sec:open}

\begin{enumerate}
\item \textbf{Critic architecture.}
  $K$ separate critics (one per axis) vs.\ one critic with $K$-headed
  output.  The headed critic shares representation and is cheaper, but
  may produce correlated advantages---exactly the degeneracy we want to
  avoid.

\item \textbf{Compute cost.}
  $K$ backward passes is $3$--$5\times$ the trainer cost at $K{=}3$--$5$.
  Is there a single-backward variant that approximates MGDA?  Candidates:
  (a) MGDA-Approx (one backward with mixed advantage, one with
  worst-axis advantage); (b) Hessian-vector-product--based MGDA that
  avoids explicit per-axis backward.

\item \textbf{Stability with GRPO.}
  GRPO's group-mean centring can push $\bm{\alpha}^*$ away from
  improving directions early in training, when per-axis advantages are
  noisy.  Does MGDA's ``stop when no improving direction exists''
  behaviour interact badly with this?

\item \textbf{Adaptive $K$.}
  Can we learn which axes to track?  Start with $K{=}10$, prune axes
  whose $\alpha_k^*$ drops below a threshold.  This is an $L_0$
  regularisation problem on the simplex.

\item \textbf{Generalisation beyond bash.}
  In math (step-level checks), code (unit-test granularity), or agentic
  tasks (intermediate state visibility), privileged teachers with
  structured feedback are more naturally available.  Does VPO generalise
  and does the axis-independence condition hold in those domains?

\item \textbf{Theoretical analysis.}
  Under what conditions does MGDA-VPO converge to a policy that
  Pareto-dominates the scalarised baseline?  Can we bound the
  improvement as a function of the inter-axis correlation matrix?
\end{enumerate}

% ======================================================================
\section{Positioning Relative to the Current Paper}
\label{sec:positioning}

The three findings of the current paper are complete and do not depend
on VPO:
\begin{enumerate}
\item Variational self-consistency + weight sync = $+15$ pp (Exp9/10).
\item Scalar self-similarity is a surprisingly strong baseline ($+13$ pp;
  the variational machinery's incremental contribution is small).
\item DBPO (density reward) is a documented failure mode (format collapse).
\end{enumerate}

VPO is naturally a \textbf{follow-up paper} about \emph{how} to design
the reward axes and \emph{how} to combine them, building on the same
weight-synced GRPO + bash environment substrate.  The current paper
establishes the substrate; VPO asks the next question.

% ======================================================================
\section*{References}
\addcontentsline{toc}{section}{References}

\begin{itemize}[nosep]
\item D\'{e}sid\'{e}ri, J.-A. (2012).
  \emph{Multiple-gradient descent algorithm for multiobjective
  optimization}.  Comptes Rendus Mathematique, 350(5--6), 313--318.
\item Sener, O. \& Koltun, V. (2018).
  \emph{Multi-task learning as multi-objective optimization}.
  NeurIPS 2018.
\item Ng, A.~Y., Harada, D. \& Russell, S. (1999).
  \emph{Policy invariance under reward transformations: Theory and
  application to reward shaping}.  ICML 1999.
\item Wang, T. \& Isola, P. (2020).
  \emph{Understanding contrastive representation learning through
  alignment and uniformity on the hypersphere}.  ICML 2020.
\end{itemize}

\end{document}
